\section{Wach lattices from one prismatic realization}
\label{sec:wach-lattice-descent-v053}

\status{This chapter proves the algebraic descent, local-prism comparison,
and lattice assembly used in the original Wach-lattice producer. The
lattice is a submodule of the \emph{supplied} \'etale module, with its
original Frobenius, group action and maps of lattices. The final all-layer
return retains three precisely stated period-ring suppliers: the Wach
prism, injectivity and reflection of units for the coefficient embedding,
and descent of cyclotomic-kernel-fixed Frobenius matrix solutions.
Their existence is not inferred from the assembly proved here.}

The obstruction addressed here is more concrete than the instruction to
choose a coherent Wach lattice suggests. A crystalline lattice has
a prismatic realization. Its value on one local prism gives a lattice
inside a Witt-vector tensor module. We first show that this lattice
comes from the given \'etale module, then take its inverse image there.
Fullness of the crystalline realization transports every lattice
homomorphism through this construction. Thus the same inverse-image
definition provides all the required functoriality.

The finite-dimensional coefficient arguments are written out because
neither faithful flatness of an unspecified period-ring extension nor
a previously constructed Wach functor is an input. In particular,
the crucial base-change isomorphism below follows from a determinant
argument and a tensor normal form. All rings in the algebraic lemmas
are commutative and all sums over finite sets include the empty case.

\subsection{Coefficients, completion and denominator matrices}

% BEGIN WLD1
\paragraph{WLD1: tensor cancellation with its actual maps.}
\label{claim:WLD1}
Let $Z\to O$ and $Z\to W$ be ring maps and let $T$ be an $O$-module.
Put $B=O\otimes_ZW$. There is a natural additive isomorphism
\[
 \kappa_T:W\otimes_ZT\longrightarrow B\otimes_OT,\qquad
 w\otimes t\longmapsto(1\otimes w)\otimes t.                 \tag{WLD.1}
\]
With $c(w)=1\otimes w$, it satisfies
$\kappa_T(wy)=c(w)\kappa_T(y)$. It commutes with every map of
lattices and every coefficient operator whose actions agree on
$c(W)$. If $b_1,\ldots,b_q$ is a $W$-basis of $W\otimes_ZT$,
then every $y\in B\otimes_OT$ has a unique expression
\[
                 y=\sum_i c(w_i)\kappa_T(b_i),\qquad w_i\in W.
                                                               \tag{WLD.2}
\]

\begin{proof}
Define the inverse on pure tensors by
\[
             (o\otimes w)\otimes t\longmapsto w\otimes ot.
\]
The $Z$-balancing relation in the first tensor product and the
$O$-balancing relation in the second are both preserved, so this
defines an additive map. Its two composites with (WLD.1) are the
identity on pure tensors and hence on all tensors. The scalar
formula follows on pure tensors from multiplication in $W$.
For operators $\sigma:W\to W$, $\sigma':B\to B$ with
$\sigma'c=c\sigma$ and maps $f,f'$ agreeing on $T$, both possible
composites send $w\otimes t$ to $c(\sigma w)\otimes f(t)$.
Additivity proves the assertion for arbitrary tensors. Finally apply
$\kappa_T$ to the unique basis expansion of $\kappa_T^{-1}(y)$.
These checks apply to the Frobenius and Galois operators as well as
to homomorphisms of the original integral lattices.
\end{proof}
% END WLD1

% BEGIN WLD2
\paragraph{WLD2: finite coefficient coordinates and reflection of units.}
\label{claim:WLD2}
If $o_1,\ldots,o_d$ is a $Z$-basis of $O$, every element of
$O\otimes_ZW$ is uniquely $\sum_j o_jc(w_j)$. If this tensor ring
is nonzero, $c(w)$ being a unit implies that $w$ is a unit of $W$.
If $T$ is free of rank $r$ over $O$, its underlying $Z$-module and
$W\otimes_ZT$ are free of rank $dr$ over $Z$ and $W$, respectively.

\begin{proof}
Tensor the coordinate isomorphism $O\simeq Z^d$ with $W$.
Multiplication by $c(w)$ is multiplication by $w$ on each of the $d$
coordinates. If it is invertible, its determinant $w^d$ is a unit.
The nonzero tensor ring excludes $d=0$; a positive power is a unit
only if its base is, since $w\cdot w^{d-1}(w^d)^{-1}=1$.
For the last assertion, if $t_1,\ldots,t_r$ is an $O$-basis of $T$,
the products $o_jt_k$ form a $Z$-basis. Tensoring gives the stated
$W$-basis. This also covers $r=0$.
\end{proof}
% END WLD2

% BEGIN WLD3
\paragraph{WLD3: lifting coefficient decompositions through completion.}
\label{claim:WLD3}
Let $O,E,W,B$ be rings, let $\iota:E\to B$ be an $O$-algebra map
and $c:W\to B$ a ring map. Suppose $o_1,\ldots,o_d\in O$ give
unique coordinates
\[
                       B=\left\{\sum_j o_jc(w_j):w_j\in W\right\}.
\]
Assume $E$ and $W$ are separated and complete for the ideals generated
by an integer $p$. If every $e\in E$ has an expression modulo $pE$
\[
 e\equiv\sum_j o_je_j\pmod {pE},\qquad \iota(e_j)\in c(W),
\]
then it has an exact expression of the same form:
\[
                  e=\sum_j o_je_j,\qquad \iota(e_j)\in c(W).
                                                               \tag{WLD.3}
\]
No division by $p$, injectivity of $c$, or a topology on $B$ is assumed.

\begin{proof}
Choose, for every $u\in E$, elements $g_j(u)\in E$, $w_j(u)\in W$
and $h(u)\in E$ such that
\[
 \iota(g_j(u))=c(w_j(u)),\qquad
                  u=\sum_j o_jg_j(u)+ph(u).
\]
Starting with $u_0=e$, put $u_{k+1}=h(u_k)$. Iteration gives
\[
 e=\sum_j o_j\sum_{k<N}p^kg_j(u_k)+p^Nu_N.                 \tag{WLD.4}
\]
Completeness supplies sums
$e_j=\sum_{k\ge0}p^kg_j(u_k)$ in $E$ and
$v_j=\sum_{k\ge0}p^kw_j(u_k)$ in $W$: concretely, the difference
from the $N$th partial sum is divisible by $p^N$.

The coefficient hypothesis makes $B$ separated for divisibility by
$p$. Indeed, if $b\in\bigcap_Np^NB$, express $b$ and each witness
$b=p^Nb_N$ in the unique $W$-coordinates. Uniqueness shows each
coordinate of $b$ is divisible by every $p^N$, hence is zero by
separatedness of $W$. Now
$\iota(e_j)-c(v_j)$ is divisible by every $p^N$ in $B$, because its
partial sums agree. Therefore it is zero. Similarly (WLD.4) shows
$e-\sum_j o_je_j$ belongs to every $p^NE$, so it is zero.
This proves (WLD.3), including cases where $p$ is a zero divisor.
\end{proof}
% END WLD3

% BEGIN WLD4
\paragraph{WLD4: coefficient decompositions in the supplied period ring.}
\label{claim:WLD4}
Let $A\hookrightarrow E$ be the power-series coefficient ring and its
localization-completion in the original \texttt{WSer} data. Write
$\pi_A$ for \texttt{letpi}. Let $a:S\to A$, $\psi:S\to W$ and
$c:W\to B$ satisfy
\[
       \iota(a(t))=c(\psi(t)),\qquad a(\pi)=\pi_A,
                                                               \tag{WLD.5}
\]
where $A\to E$ is suppressed in notation. Suppose every $\alpha\in A$
is $\sum_j o_ja(t_j)$, $B$ has the unique $W$-coordinates of WLD3,
$W$ is $p$-adically complete, and $c$ reflects units. Then the
original ring $E$ satisfies (WLD.3).

\begin{proof}
The original density assertion, applied with precision one, gives
$\alpha\in A$ and $m\ge0$ such that
$e\pi_A^m-\alpha\in pE$. The element $\pi_A$ is a unit in $E$.
Its image is $c(\psi(\pi))$, so $\psi(\pi)$ is a unit in $W$.
Write $\alpha=\sum_j o_ja(t_j)$ and set
$e_j=\pi_A^{-m}a(t_j)$. Then
\[
 e-\sum_j o_je_j\in pE,\qquad
 \iota(e_j)=c\bigl(\psi(\pi)^{-m}\psi(t_j)\bigr).
\]
The first equality uses multiplication by the unit $\pi_A^{-m}$;
it does not cancel $p$. Apply WLD3, using the completeness of $E$
included in \texttt{WSer}. Thus the desired exact coefficients
are obtained in the \emph{same} $E$, rather than by replacing it with
a new completion.
\end{proof}
% END WLD4

% BEGIN WLD5
\paragraph{WLD5: independence after clearing one common denominator.}
\label{claim:WLD5}
Let $a:S\to E$, $\iota:E\hookrightarrow B$, $c:W\to B$,
$\psi:S\to W$ satisfy $\iota a=c\psi$, and let $a(Q)$ be a unit.
Let $y_1,\ldots,y_q$ in a $B$-module be independent for coefficients
from $W$, meaning
\[
             \sum_i c(w_i)y_i=0\ \Longrightarrow\ w_i=0.
\]
Define the subring
\[
 E_Q=\{e\in E:\ a(Q)^ne=a(t)
                    \text{ for some }n\ge0,\ t\in S\}.
\]
Then $\sum_i\iota(e_i)y_i=0$ with $e_i\in E_Q$ implies $e_i=0$.

\begin{proof}
The set $E_Q$ is closed under sums, products and negatives by putting
fractions over a common power of $a(Q)$. Choose presentations
$a(Q)^{n_i}e_i=a(t_i)$ and take $N\ge n_i$ for every $i$.
Multiplying the relation by $\iota(a(Q)^N)$ gives
\[
        \sum_i c\bigl(\psi(Q^{N-n_i}t_i)\bigr)y_i=0.
\]
Independence makes every displayed $W$-coefficient zero.
The identity $\iota a=c\psi$ and injectivity of $\iota$ imply
$a(Q)^Ne_i=0$. Since $a(Q)$ is a unit, $e_i=0$.
\end{proof}
% END WLD5

% BEGIN WLD6
\paragraph{WLD6: two Frobenius span inclusions give an invertible matrix.}
\label{claim:WLD6}
In WLD5, suppose another tuple $z_1,\ldots,z_q$ satisfies
\[
 c(\psi Q)^{k_j}z_j=\sum_i c(\psi(s_{ij}))y_i,\qquad
 c(\psi Q)^{h_i}y_i=\sum_j c(\psi(t_{ji}))z_j
\]
for suitable nonnegative exponents and coefficients in $S$.
Then there is $U\in\operatorname{GL}_q(E)$ with
\[
                         z_j=\sum_i\iota(U_{ij})y_i.       \tag{WLD.6}
\]

\begin{proof}
Invert $a(Q)$ in $E$ in the first family of equations. It gives a
matrix $U$ with entries in $E_Q$ and the displayed expression for
$z_j$. The second family similarly gives a matrix $V$ over $E_Q$
expressing the $y_i$ in terms of the $z_j$. Substitute to obtain
$\sum_m\iota((UV-I)_{mi})y_m=0$. WLD5 gives $UV=I$.
For square matrices over a commutative ring this implies
$\det(U)\det(V)=1$; the adjugate formula then makes $U$ invertible.
Neither the ambient module nor the coefficient extension is assumed
flat in this argument.
\end{proof}
% END WLD6

\subsection{The exact base-change and projectivity arguments}

% BEGIN WLD7
\paragraph{WLD7: spanning the supplied module by a descended tuple.}
\label{claim:WLD7}
Let $O\to A\to E$, let $\iota:E\hookrightarrow B$ be an $O$-algebra
map reflecting units, and let $a:S\to A$, $c:W\to B$ be ring maps.
Let $D$ be free of rank $r$ over $E$, $Y$ a $B$-module, and
$j:D\to Y$ additive with $j(ex)=\iota(e)j(x)$. Assume
$\iota(a(S))\subseteq c(W)$ and
\[
 \forall e\in E\quad
 e=\sum_{j=1}^d o_je_j\quad\text{with }\iota(e_j)\in c(W).
                                                               \tag{WLD.7}
\]
Suppose $q=dr$ and $x_1,\ldots,x_q\in D$ give unique $W$-coordinates
on $Y$ through their images:
\[
                    y=\sum_i c(w_i)j(x_i).
                                                               \tag{WLD.8}
\]
Then $x_1,\ldots,x_q$ span $D$ over $E$.

\begin{proof}
Choose an $E$-basis $b_1,\ldots,b_r$. Index a tuple of length $dr$
by pairs and put $u_{jk}=o_jb_k$. Decompose each coordinate of
each $x_i$ using (WLD.7). This produces a square matrix $C'$ over
$E$ and a square matrix $C$ over $W$ with
\[
             x_i=\sum_l C'_{il}u_l,\qquad
                         \iota(C'_{il})=c(C_{il}).
\]
Use (WLD.8) to express $j(u_l)=\sum_m c(B_{lm})j(x_m)$.
Substitution and uniqueness in (WLD.8) give $CB=I$.
Thus $\det C$ is a unit. Since
$\iota(\det C')=c(\det C)$ and $\iota$ reflects units,
$\det C'$ is a unit in $E$. The inverse matrix expresses every
$u_l$ as an $E$-linear combination of the $x_i$.

Finally decompose $1=\sum_j o_je_j$ using (WLD.7). For every $k$,
$b_k=\sum_j e_ju_{jk}$, so the $b_k$, and hence all of $D$, lie
in the span of the $x_i$. This proof does not claim that the
$q=dr$ vectors are an $E$-basis; generally they are not.
\end{proof}
% END WLD7

% BEGIN WLD8
\paragraph{WLD8: the tensor map is an isomorphism without an extra flatness premise.}
\label{claim:WLD8}
Keep the assumptions of WLD7. Let $N\subset D$ be an $A$-submodule
whose elements are precisely
\[
                        \sum_i a(s_i)x_i,\qquad s_i\in S.
                                                               \tag{WLD.9}
\]
Then the actual inclusion induces an isomorphism
\[
                         E\otimes_AN\xrightarrow{\ \sim\ }D.
                                                               \tag{WLD.10}
\]

\begin{proof}
Each $x_i$ lies in $N$, so WLD7 gives surjectivity.
For injectivity, first put every tensor into the normal form
\[
              t=\sum_i e_i\otimes x_i,\qquad
                                      \iota(e_i)\in c(W).
                                                               \tag{WLD.11}
\]
For a pure tensor $e\otimes n$, decompose $e=\sum_j o_je_j$.
Move $o_j\in A$ across the tensor sign. Since $N$ is an
$A$-module, $o_jn$ has an expression (WLD.9). Move its coefficients
$a(s_i)$ back across the tensor sign. The resulting coefficients
are sums of products of $e_j$ with $a(s_i)$; their images lie in
the subring $c(W)$. Additivity proves (WLD.11) for all tensors.

If (WLD.11) maps to zero, choose $w_i\in W$ with
$c(w_i)=\iota(e_i)$. Applying $j$ gives
$\sum_i c(w_i)j(x_i)=0$. Uniqueness of (WLD.8), compared with
the zero tuple, gives $w_i=0$. Injectivity of $\iota$ gives
$e_i=0$, whence $t=0$. The map in (WLD.10) is therefore bijective
and is the original base-change map on every $1\otimes n$.
This proves the full source statement
\texttt{is\allowbreak Base\allowbreak Change\_\allowbreak lattice\_\allowbreak of\_\allowbreak basis\_\allowbreak after\_\allowbreak extension}.
\end{proof}
% END WLD8

% BEGIN WLD9
\paragraph{WLD9: projectivity over the original power-series ring.}
\label{claim:WLD9}
Let $Z\to S\xrightarrow{a}A$ be ring maps, and let $N$ be an
$A$-submodule of an $A$-module. Suppose (WLD.9) holds and the
coefficients $s_i$ in that expression are unique. If
$S$ is finite free over $Z$ with basis $v_1,\ldots,v_h$, then
$N$, with the restricted $Z$-action, is free on
\[
                              a(v_j)x_i.
\]
In particular it is projective. In the Wach application
$Z=\mathbb Z_p[[X]]$, with its specified map into $S$.

\begin{proof}
Write each $s_i$ uniquely as $\sum_j z_{ij}v_j$.
This gives existence of the indicated expression for every element
of $N$. If that expression is zero, uniqueness of the $S$-coordinates
first gives $\sum_jz_{ij}v_j=0$ for every $i$; the $Z$-basis of $S$
then gives $z_{ij}=0$. The resulting explicit isomorphism
$Z^{qh}\simeq N$ proves freeness and projectivity with exactly the
restricted scalar action. We do not infer projectivity over $A$
from this argument.
\end{proof}
% END WLD9

\subsection{What a local prism gives inside the Witt tensor module}

The word \emph{crystal} below means the precise source data
\texttt{FCrys}: a finite projective module $M(P)$ on each prism,
semilinear transition maps inducing base-change isomorphisms, their
identity and composition laws, and a Frobenius isomorphism after
localization at the prism ideal. We use its stated fraction-naturality
law, so no injectivity of a transition map is silently assumed.
A realization identifies the base change at the distinguished
$A_{\inf}$ prism with $W\otimes_{\mathbb Z_p}T$, respecting Frobenius
fractions and Galois action. Here $W=W(C^\flat)$ is the source's
\texttt{WCb}; it is a $p$-adically complete domain. These are the
original period rings, not substitute abstract realizations.

% BEGIN WLD10
\paragraph{WLD10: bases and principal kernels under semilinear base change.}
\label{claim:WLD10}
Let $f:S\to R$ be a ring map and $u:M\to N$ an $f$-semilinear map
inducing an isomorphism $R\otimes_{S,f}M\simeq N$. If $M$ has
basis $b_1,\ldots,b_q$, then $u(b_1),\ldots,u(b_q)$ is an
$R$-basis of $N$. If $\ker f=(\pi)$, then
\[
                            \ker u=\pi M.                \tag{WLD.12}
\]
These conclusions require no flatness of $R$ over $S$.

\begin{proof}
Tensor the explicit coordinate isomorphism $M\simeq S^q$ with $R$
and compose it with the given base-change isomorphism.
The resulting map sends the standard basis to $u(b_i)$, proving
the first assertion. If $u(x)=0$ and $x=\sum_i s_ib_i$,
the independence of $u(b_i)$ gives $f(s_i)=0$. Write
$s_i=\pi t_i$; then $x=\pi\sum_i t_ib_i$. Conversely
$f(\pi)=0$ shows $\pi M\subseteq\ker u$.
The source's \texttt{Semi\allowbreak BC} is the equivalent universal property:
the inverse to the displayed coordinate map is the unique linear
map induced by $x\mapsto(f(s_i(x)))_i$. Thus the proof applies
directly to that formulation.
\end{proof}
% END WLD10

% BEGIN WLD11
\paragraph{WLD11: the two Frobenius inclusions at a local prism.}
\label{claim:WLD11}
Fix a local prism $P_S$ with coefficient ring $S$, an element
$Q$ of its prism ideal, and a map $i:P_S\to P_{\inf}$.
Let a crystalline lattice $T$ have a realizing crystal $\mathcal C$.
Write
\[
 M=\mathcal C(P_S),\quad
 Y_0=W\otimes_{\mathbb Z_p}T,\quad
 \nu_0=j\circ\mathcal C(i):M\longrightarrow Y_0,
\]
and let $\psi:S\to W$ be its scalar map. Suppose $\psi(Q)\ne0$.
There are bases $b_1,\ldots,b_q$ of $M$ over $S$ and
$c_1,\ldots,c_q$ of $Y_0$ over $W$, with $c_i=\nu_0(b_i)$.
For the actual Frobenius $F$ on $Y_0$ one has
\[
 \begin{split}
 &\forall m\in M\quad
   \exists k\ge0,\ m'\in M:\quad
           \psi(Q)^kF(\nu_0m)=\nu_0m',\\
 &\forall m\in M\quad
   \exists k\ge0,\ s_1,\ldots,s_q\in S:\quad
           \psi(Q)^k\nu_0m=\sum_i\psi(s_i)F(c_i).
 \end{split}                                                \tag{WLD.13}
\]

\begin{proof}
A finite projective module over a local ring is finite free.
For completeness, choose lifts of a residue-field basis, use
Nakayama to obtain a surjection from a finite free module, and split
it by projectivity. Its complementary summand has zero reduction
and is finite, hence vanishes by Nakayama. Choose the resulting
basis $b_i$. Apply WLD10 to the crystal transition and to the
realization. Their composite carries $b_i$ to a $W$-basis $c_i$.

We spell out the fraction comparison. Write the localized Frobenius
of $m$ at $P_S$ as $m'/Q^k$, and the corresponding Frobenius at
$P_{\inf}$ as $n/i(Q)^h$. Fraction naturality says that, after
multiplication by a power of $i(Q)^2$,
\[
               i(Q)^h\,\mathcal C(i)(m')
                       -i(Q)^k n
\]
is zero. Apply $j$. The realization identifies
$j(n)=\psi(Q)^hF(\nu_0m)$. Consequently a power of the nonzero
element $\psi(Q)$ annihilates
$\nu_0m'-\psi(Q)^kF(\nu_0m)$. The ambient module is free over the
domain $W$, so this difference is zero. This proves the first
inclusion without cancelling an element inside the possibly
non-domain ring $S$.

For the other inclusion, write the localized Frobenius of $b_i$
as $m_i/Q^{k_i}$. Its semilinear base-change property makes those
fractions a basis of $M[1/Q]$, by WLD10. Thus the localized class
of any $m$ lies in the span of the classes of the $m_i$.
Clear the finitely many coefficient denominators. If the resulting
equality initially holds only in the localization, multiply by
one further power of $Q$ killing its discrepancy in $M$.
We obtain an actual equality
\[
                         Q^k m=\sum_i t_i m_i
\]
in $M$. Applying $\nu_0$ and using the first part for each $b_i$
gives
\[
 \psi(Q)^k\nu_0m
       =\sum_i\psi(t_iQ^{k_i})F(c_i),
\]
which is the second assertion. The case $q=0$ gives zero modules
and both assertions with exponent zero.
\end{proof}
% END WLD11

% BEGIN WLD12
\paragraph{WLD12: group action and divisibility at the specialization.}
\label{claim:WLD12}
In WLD11, suppose the source supplies maps
$g_S:P_S\to P_S$ and $e:P_S\to P_W$ such that
\[
 i\circ g_S=g_{\inf}\circ i,\qquad e\circ g_S=e,\qquad
                         \ker(e:S\to P_W.A)=(\pi).
\]
Then the realized action satisfies
\[
 g(\nu_0m)=\nu_0(\mathcal C(g_S)m),\qquad
                 \mathcal C(g_S)m-m=\pi m_g
                         \quad\text{for some }m_g\in M.    \tag{WLD.14}
\]
Thus WLD11--WLD12 give every conclusion of the original
\texttt{prismatic\allowbreak FCrystal\_\allowbreak value\_\allowbreak local\allowbreak Prism\_\allowbreak ainf\_\allowbreak witt\allowbreak Tensor}
on its supplied local prism and realization.

\begin{proof}
The first equality follows by applying the transition composition
law to $i\circ g_S=g_{\inf}\circ i$, then the Galois compatibility
of $j$. For the second, transition along $e$ annihilates
$\mathcal C(g_S)m-m$ because $e\circ g_S=e$.
That transition has the semilinear base-change property.
Apply (WLD.12) with the basis of $M$ from WLD11.
This is a divisibility statement in the original $M$, with the
original $\pi$; it is stronger than mere equality after localization.
\end{proof}
% END WLD12

% BEGIN WLD13
\paragraph{WLD13: preservation of realized images under every lattice map.}
\label{claim:WLD13}
Assume the original \texttt{Bhatt\allowbreak Scholze\allowbreak Crystalline\allowbreak Lattices}
statement. Choose any realizing crystals for crystalline lattices
$T,T'$. If $f:T\to T'$ is a $\mathbb Z_p$-linear equivariant map,
then
\[
 (1\otimes f)\bigl(\nu_{0,T}(M_T)\bigr)
                          \subseteq \nu_{0,T'}(M_{T'}).    \tag{WLD.15}
\]
In particular this holds for every original $O$-linear lattice
homomorphism, and for multiplication by every $o\in O$ on $T$.

\begin{proof}
The full-faithfulness clause of the stated crystalline-lattice input
gives a morphism $u:\mathcal C_T\to\mathcal C_{T'}$ realizing $f$.
At the local prism use $u_{P_S}(m)$. Naturality along $i$ and the
identity defining the realization of $u$ give
\[
 \nu_{0,T'}(u_{P_S}(m))
       =(1\otimes f)\nu_{0,T}(m).
\]
An $O$-linear lattice map restricts to a $\mathbb Z_p$-linear map
and remains equivariant, since its underlying function is unchanged.
The same is true for scalar multiplication by $o$.
The proof applies to the crystals already chosen for $T$ and $T'$;
it does not require replacing either crystal by a newly compatible
choice. The uniqueness clause is available, but existence is enough
for this inclusion.
\end{proof}
% END WLD13

% BEGIN WLD14
\paragraph{WLD14: the exact image for a trivial lattice.}
\label{claim:WLD14}
Under WLD13, let $T$ have trivial action and a
$\mathbb Z_p$-basis $\tau_1,\ldots,\tau_d$. Then
\[
 \nu_{0,T}(M_T)
       =\left\{\sum_j\psi(s_j)\otimes\tau_j:s_j\in S\right\}.
                                                               \tag{WLD.16}
\]
For the source's specified trivial rank-one $O$-lattice,
choose its generator $e_0$ and a $\mathbb Z_p$-basis $o_j$ of $O$.
If every $\alpha\in A$ is $\sum_j o_ja(s_j)$, WLD1 carries
(WLD.16) to the exact subset
\[
      \left\{\iota(\alpha)(1\otimes e_0):\alpha\in A\right\}
                   \subseteq B\otimes_OT.                 \tag{WLD.17}
\]

\begin{proof}
The unit crystal has value the coefficient ring on each prism,
transition the ring map, and localized Frobenius the Frobenius
ring map. These maps have the required semilinear base-change
property because the module has basis $1$. Its realization sends
$s\in S$ to $\psi(s)\otimes1$.

Each coordinate projection $T\to\mathbb Z_p$ is equivariant.
Apply WLD13 to it and to the unit crystal. Hence the corresponding
coordinate of any $\nu_{0,T}(m)$ lies in $\psi(S)$. Reassembling
the coordinates proves inclusion in the right-hand side of
(WLD.16). Conversely the map $1\mapsto\tau_j$ is equivariant;
WLD13 applied in this direction shows
$\psi(s_j)\otimes\tau_j$ belongs to the image. Additivity proves
the reverse inclusion.

For the specified rank-one lattice, its nonzero line generator and
injectivity $O\hookrightarrow LL$ identify its integral lattice
with $Oe_0$. Consequently $\tau_j=o_je_0$ is a
$\mathbb Z_p$-basis. Under WLD1,
\[
 \sum_j\psi(s_j)\otimes o_je_0
       \longmapsto
       \left(\sum_j o_jc(\psi(s_j))\right)(1\otimes e_0)
       =\iota\left(\sum_j o_ja(s_j)\right)(1\otimes e_0).
\]
The assumed coefficient decomposition of $A$ proves (WLD.17)
in both directions. This fixes the actual trivial lattice, not
only its isomorphism class.
\end{proof}
% END WLD14

\subsection{Returning the prismatic basis to the given \'etale module}

Fix the original group, series, embedding and \'etale data
$\Gamma,R,I,\mathcal D$. In ordinary notation put
\[
 A=\operatorname{LetA}(O,n),\quad E=R.\mathrm{EE},\quad
 B=R_W,\quad D_T=\mathcal D(T),\quad Y_T=B\otimes_OT.
\]
The map $j_T:D_T\to Y_T$ is injective and
$j_T(ex)=\iota(e)j_T(x)$. Its image spans $Y_T$ over $B$.
The functor has its supplied semilinear Frobenius $\varphi_D$,
$\Gamma$-action $\gamma_D$, and maps $D(f)$; $j_T$ intertwines
these with the actual tensor operators. The modules $D_T$ are
free and have the stated rank of $T$. The maps $D(f)$ already
satisfy identity, composition, addition and scalar laws.
These are precisely the original \texttt{WEt} fields and are
inputs here, not new existence assumptions.

% BEGIN WLD15
\paragraph{WLD15: fixed Frobenius tuples lie in the given \'etale module.}
\label{claim:WLD15}
Suppose $\iota$ reflects units, and $H\subset G$ maps trivially to
$\Gamma$. Fix $x_1,\ldots,x_q\in Y_T$ and a matrix $U\in M_q(E)$
with
\[
 h(x_i)=x_i\ (h\in H),\qquad
            \varphi_Y(x_j)=\sum_i\iota(U_{ij})x_i.          \tag{WLD.18}
\]
Assume the following precise coefficient-descent assertion for this
$U$: for every $r$, $F\in\operatorname{GL}_r(E)$ and
$V\in M_{r,q}(B)$, if its entries are $H$-fixed and
\[
                    \iota(F)\varphi_B(V)=V\iota(U),        \tag{WLD.19}
\]
then all its entries belong to $\iota(E)$. Then every $x_i$ lies
in $j_T(D_T)$.

\begin{proof}
If $B=0$ or $T$ has rank zero, $Y_T=0$, and zero supplies the
required preimage. Otherwise the original rank and freeness
fields give a finite $E$-basis $b_1,\ldots,b_r$ of $D_T$.
The $j_T(b_i)$ span $Y_T$, which is free of rank $r$ over $B$:
expand every $j_T(d)$ in those vectors, and use the spanning
field of $j_T$. A spanning $r$-tuple in a free module of rank $r$
is a basis. Indeed, its square coordinate matrix is surjective,
has a right inverse, and hence has invertible determinant.

Let $F$ be the matrix of $\varphi_D$ on this $E$-basis. Its
image in $B$ is the semilinear Frobenius matrix in the
$j_T(b_i)$ basis. That matrix is invertible. To see this without
assuming it, note that the fixed vectors $1\otimes t$ span $Y_T$.
For a fixed vector with coordinate column $v$ we have
$v=\iota(F)\varphi_B(v)$. Thus the columns of $\iota(F)$ span
$B^r$ and its determinant is a unit. Reflection of units makes
$\det F$ a unit in $E$.

Express $x_j=\sum_iV_{ij}j_T(b_i)$. Since $H$ acts trivially
on $\Gamma$, it fixes every $j_T(b_i)$; uniqueness of coordinates
and the first part of (WLD.18) make every $V_{ij}$ $H$-fixed.
The second part of (WLD.18), expanded in this basis, is exactly
(WLD.19). Choose $v_{ij}\in E$ with $\iota(v_{ij})=V_{ij}$.
Then $x_j=j_T(\sum_i v_{ij}b_i)$.

There is no hidden finite-rank inference when the rank field is
zero: the finite free module $Y_T$ is then zero, and injectivity
of $j_T$ forces $D_T=0$. This also covers the source convention
that natural-number finrank alone need not establish finite
generation of an arbitrary free module.
\end{proof}
% END WLD15

% BEGIN WLD16
\paragraph{WLD16: the supplied trivial object is crystalline.}
\label{claim:WLD16}
For the original period datum, every finite-dimensional trivial
representation is crystalline. In particular this applies to the
specified trivial line $\Gamma.\mathrm{one}$, including after any
allowed finite extension of coefficient fields.

\begin{proof}
Let $V$ be its underlying finite-dimensional $\mathbb Q_p$-space.
Choose a basis $v_i$. Write a tensor in
$V\otimes_{\mathbb Q_p}B_{\mathrm{cris}}$ uniquely as
$\sum_i v_i\otimes b_i$. Since the action on $V$ is trivial,
this tensor is invariant exactly when each $b_i$ is invariant.
The invariant-ring identity of the original period datum is
$B_{\mathrm{cris}}^G=K_0$. Hence
\[
            D_{\mathrm{cris}}(V)
              =V\otimes_{\mathbb Q_p}K_0
              \subseteq V\otimes_{\mathbb Q_p}B_{\mathrm{cris}}.
\]
The inclusion is injective because $K_0$ is a field and its map to
the nonzero period ring is injective. Taking dimensions gives
precisely the original crystalline dimension equality.
Finite dimension over $\mathbb Q_p$ follows from finite dimension
over the allowed finite coefficient field. This proves the entire
\texttt{is\allowbreak Crys\_\allowbreak lat\_\allowbreak of\_\allowbreak trivial} assertion on the supplied period data.
\end{proof}
% END WLD16

% BEGIN WLD17
\paragraph{WLD17: one prismatic basis inside the actual \'etale module.}
\label{claim:WLD17}
Assume also that $\iota:E\to B$ is injective, as in the original
\texttt{build\_prism} hypotheses.
Keep one local prism and all its maps as in WLD11--WLD14. Suppose
$a:S\hookrightarrow A$ and $\psi:S\to W$ satisfy (WLD.5),
$a(Q)=\mathcal Q$, and the image of $\mathcal Q$ is a unit in $E$.
Assume the source's fixed-tuple descent property: every $H$-fixed
finite tuple in $Y_T$ whose Frobenius is given by an invertible
matrix over $E$ belongs to $j_T(D_T)$. Here $H$ maps trivially
to $\Gamma$, and $a(g_S(t))=\gamma_A(\operatorname{res}_\Gamma g)a(t)$.
Then, for each crystalline $T$, one can choose $x_1,\ldots,x_q\in D_T$
with $q=dr_T$, where $d$ is the $\mathbb Z_p$-rank of $O$ and
$r_T$ the $O$-rank of $T$, such that
\[
 \begin{split}
 &\text{every }y\in Y_T
          \text{ is uniquely }\sum_i c(w_i)j_T(x_i),\\
 &L_T:=\kappa_T(\nu_{0,T}(M_T))
       =\left\{\sum_i c(\psi(s_i))j_T(x_i):s_i\in S\right\}.
 \end{split}                                                \tag{WLD.20}
\]
The actual fixed-tuple descent can be supplied by WLD15; it is not
a consequence of nonvanishing or of equality of ranks alone.

\begin{proof}
Take the bases $b_i,c_i$ of WLD11 and set
$y_i=\kappa_T(c_i)$. WLD1 gives their unique $W$-coordinates on
$Y_T$. Since $T$ is free of rank $r_T$ over $O$ and $O$ free
of rank $d$ over $\mathbb Z_p$, WLD2 and invariance of finite free
rank over the nonzero ring $W$ give $q=dr_T$.

The two Frobenius inclusions (WLD.13), transported by WLD1,
give precisely the hypotheses of WLD6 for the tuples
$y_i$ and $\varphi_Y(y_i)$. Thus the Frobenius matrix on this
tuple is invertible over $E$.

If $g\in H$, equivariance of $a$ gives $a(g_S(t))=a(t)$
for every $t\in S$. Its injectivity implies that $g_S$ is the
identity prism map. WLD12 and the identity law for crystal
transition maps then give $g(y_i)=y_i$. The supplied fixed-tuple
descent provides $x_i$ with $j_T(x_i)=y_i$. These preimages are
unique because $j_T$ is injective. Finally every element of $M_T$
has unique $S$-coordinates in the basis $b_i$; semilinearity
of $\nu_{0,T}$ and WLD1 give the second line of (WLD.20).
The chosen basis helps prove the statement, but the subset $L_T$
is the image of the chosen prismatic realization itself.
\end{proof}
% END WLD17

\subsection{Constructing and checking every field of the lattice}

We now prove the original lattice assembly at its full generic scope.
This also isolates exactly what the prismatic realization supplies.
Keep $A,E,B,D_T,j_T$ and all their original operators fixed.
Suppose $E\ne0$ and the following data have been supplied:
\begin{enumerate}
 \item Ring maps $Z_0[[X]]\xrightarrow{z}S\xrightarrow{a}A$,
 with $a$ injective and $S$ finite free over $Z_0[[X]]$;
 maps $c:W\to B$, $\psi:S\to W$ satisfying (WLD.5);
 $\iota:E\hookrightarrow B$ reflecting units;
 and elements $\pi,Q\in S$ with
 $a(\pi)=\pi_A$ and $a(Q)=\mathcal Q$.
 \item Elements $o_1,\ldots,o_d\in O$ such that
 $A=\sum_j o_ja(S)$, and the exact coefficient decomposition
 (WLD.7) in $E$.
 \item For every crystalline $T$, a subset $L_T\subseteq Y_T$
 and a tuple $x_i$ satisfying both assertions (WLD.20), with
 $q=dr_T$.
 \item For every crystalline $T$ and $y\in L_T$, some power
 of $c(\psi Q)$ times $\varphi_Y(y)$ lies in $L_T$, and some
 power times $y$ is a finite sum of terms
 $c(\psi t)\varphi_Y(u)$ with $t\in S$, $u\in L_T$.
 For every $g\in G$ and $y\in L_T$ there is $u\in L_T$ with
 $g(y)-y=c(\psi\pi)u$.
 \item The subset $L_T$ is preserved by every tensor map
 associated to an original homomorphism between crystalline lattices.
 \item The specified trivial object $T_1$ is crystalline, has an
 actual generator $e_0$ over $O$, and
 $L_{T_1}=\{\iota(\alpha)(1\otimes e_0):\alpha\in A\}$.
\end{enumerate}
These are the explicit hypotheses of the source's
\texttt{Wach\allowbreak Lattice.\allowbreak build}. They concern subsets of the ambient tensor
modules and finite coefficient data, not an already existing lattice
in $D_T$. In the prismatic application, WLD11--WLD17 prove the
subset assertions on one fixed prism, while WLD4 supplies the
coefficient decomposition.

% BEGIN WLD18
\paragraph{WLD18: the inverse image is a submodule over the full coefficient ring.}
\label{claim:WLD18}
For crystalline $T$ define
\[
           N(T)=\{x\in D_T:j_T(x)\in L_T\};
\]
for noncrystalline $T$ put $N(T)=0$. For crystalline $T$ this is
an $A$-submodule, and
\[
 x\in N(T)\quad\Longleftrightarrow\quad
                       x=\sum_i a(s_i)x_i\quad(s_i\in S).
                                                               \tag{WLD.21}
\]
Every element of $L_T$ has a unique preimage under $j_T$.

\begin{proof}
The second assertion of (WLD.20), semilinearity of $j_T$ and its
injectivity give (WLD.21). They also give existence and uniqueness
of the stated preimages. Formula (WLD.21) proves closure under
addition, zero and multiplication by $a(S)$.

For $o\in O$, apply the supplied map-preservation assertion to the
original lattice endomorphism $t\mapsto ot$. The scalar law of
the given \'etale functor and naturality of $j_T$ show that
$j_T(ox)$ is in $L_T$ whenever $j_T(x)$ is. Thus the inverse
image is stable under $O$ as well. Now write an arbitrary
$\alpha\in A$ as $\sum_j o_ja(s_j)$. The preceding two stability
statements and addition prove stability under $\alpha$.
This last step is essential: stability only under the smaller
ring $S$ would not supply the original \texttt{WN.\allowbreak N} field.
\end{proof}
% END WLD18

% BEGIN WLD19
\paragraph{WLD19: finite generation, actual scalar extension and projectivity.}
\label{claim:WLD19}
The submodule $N(T)$ is finitely generated over $A$ for every $T$.
For crystalline $T$ the inclusion gives the original isomorphism
\[
                        E\otimes_A N(T)\simeq D_T.         \tag{WLD.22}
\]
For every $T$, its restriction along $a\circ z$ is projective
over $Z_0[[X]]$.

\begin{proof}
For crystalline $T$, every $x_i$ lies in $N(T)$ by (WLD.21),
and that formula expresses every element as an $A$-linear
combination of the finite tuple. Conversely their $A$-span is
in $N(T)$ by WLD18. Thus they generate exactly $N(T)$.

To apply WLD8, one needs a finite $E$-basis of $D_T$ of rank $r_T$.
If $r_T>0$, the original freeness and finrank equality provide it.
If $r_T=0$, then $q=0$ and the unique-coordinate assertion
(WLD.20) gives $Y_T=0$; injectivity of $j_T$ gives $D_T=0$.
The empty basis suffices. WLD7--WLD8 now prove (WLD.22) on
the literal inclusion.

For projectivity, suppose $\sum_i a(s_i)x_i=0$. Applying $j_T$
gives $\sum_i c(\psi(s_i))j_T(x_i)=0$. Uniqueness of the
$W$-coordinates makes $\psi(s_i)=0$ for every $i$.
Then $\iota(a(s_i))=0$. The successive injective maps
$S\hookrightarrow A\hookrightarrow E\hookrightarrow B$
give $s_i=0$. Hence the $S$-coordinates in (WLD.21) are unique.
WLD9 proves that $N(T)$ is finite free, in particular projective,
over $Z_0[[X]]$ with precisely the specified restricted action.
For noncrystalline $T$, the zero module is finitely generated
and projective. No scalar-extension assertion is required in
that branch.
\end{proof}
% END WLD19

% BEGIN WLD20
\paragraph{WLD20: both original Frobenius localization conditions.}
\label{claim:WLD20}
For every $T$ and $x\in N(T)$ there is $k\ge0$ with
\[
                   \mathcal Q^k\varphi_D(x)\in N(T).
\]
For every $y\in N(T)$ there is $k\ge0$ with
\[
           \mathcal Q^ky\in
                  \operatorname{span}_A\{\varphi_D(x):x\in N(T)\}.
                                                               \tag{WLD.23}
\]
Here $\mathcal Q$ is the original
$\varphi_A^{\,s-1}(q_0)$, not a replacement denominator.

\begin{proof}
For crystalline $T$, $j_T$ intertwines Frobenius and carries
$\mathcal Q^kx$ to $c(\psi Q)^kj_T(x)$. The first subset
condition in item 4 therefore proves the first assertion.

For the second, write the corresponding ambient relation as
\[
 c(\psi Q)^kj_T(y)=
                 \sum_i c(\psi(t_i))\varphi_Y(u_i),
                         \qquad u_i\in L_T.
\]
By WLD18 lift each $u_i$ to $v_i\in N(T)$. Intertwining and
injectivity of $j_T$ give the actual equality
\[
                      \mathcal Q^ky=\sum_i a(t_i)\varphi_D(v_i)
\]
in $D_T$. This lies in the span in (WLD.23).
For noncrystalline $T$, $x=y=0$, and exponent zero works.
The exponents are permitted to depend on the vector, exactly
as in the original two \texttt{WN} fields.
\end{proof}
% END WLD20

% BEGIN WLD21
\paragraph{WLD21: group stability and trivial action modulo the fixed parameter.}
\label{claim:WLD21}
For every $\gamma\in\Gamma$ and $x\in N(T)$ there exists
$u\in N(T)$ with
\[
                        \gamma_D(x)-x=\pi_Au.             \tag{WLD.24}
\]
Consequently $N(T)$ is $\Gamma$-stable and its induced action
on $N(T)/\pi_A N(T)$ is trivial.

\begin{proof}
For crystalline $T$, use the surjectivity $G\to\Gamma$ in the
original group data to lift $\gamma$ to $g\in G$. Item 4 gives
$g(j_Tx)-j_Tx=c(\psi\pi)y$ for some $y\in L_T$.
Lift $y=j_Tu$ by WLD18. Compatibility of the action, (WLD.5)
and injectivity of $j_T$ give (WLD.24).
Since $N(T)$ is an $A$-submodule, the equality also gives
$\gamma_Dx=x+\pi_Au\in N(T)$.
For noncrystalline $T$ choose $u=0$. Notice that the divisibility
is inside $N(T)$, rather than merely inside the ambient $D_T$.
\end{proof}
% END WLD21

% BEGIN WLD22
\paragraph{WLD22: all original lattice homomorphisms preserve the construction.}
\label{claim:WLD22}
For any original lattice homomorphism $f:T\to T'$ with $T,T'$
crystalline, the supplied map $D(f)$ takes $N(T)$ into $N(T')$.
Its restrictions preserve the original identity, composition,
addition, scalar, Frobenius and $\Gamma$ relations.

\begin{proof}
If $x\in N(T)$, item 5 puts
$(1\otimes f)j_Tx$ in $L_{T'}$. The actual naturality identity
\[
                       j_{T'}D(f)x=(1\otimes f)j_Tx
\]
and the definition of $N(T')$ give the inclusion.
All the asserted map identities are equalities of restrictions
of the given maps on $D_T$; their images lie in the required
submodules by this inclusion and WLD20--WLD21. The original
Frobenius relation may be read after the prescribed localization,
where WLD20 makes the restrictions defined.
No new choices of matrices or bases are made for different $f$.
The original \texttt{map\_\allowbreak N} field only requires both lattices
to be crystalline, and we retain exactly that scope.
\end{proof}
% END WLD22

% BEGIN WLD23
\paragraph{WLD23: the specified trivial lattice and its generator.}
\label{claim:WLD23}
There is an element $e\in D_{T_1}$ with
\[
        j_{T_1}(e)=1\otimes e_0,\qquad N(T_1)=Ae,           \tag{WLD.25}
\]
where $e_0$ is the supplied generator of the original integral
trivial lattice.

\begin{proof}
Item 6 puts $1\otimes e_0$ in $L_{T_1}$, so WLD18 gives its
unique preimage $e$. An element $x$ belongs to $N(T_1)$ exactly
when $j_{T_1}x=\iota(\alpha)(1\otimes e_0)$ for some
$\alpha\in A$. Semilinearity and injectivity say exactly that
$x=\alpha e$. Thus (WLD.25) identifies the actual chosen
trivial object and its actual tensor generator.
\end{proof}
% END WLD23

% BEGIN WLD24
\paragraph{WLD24: the full original generic Wach-lattice constructor.}
\label{claim:WLD24}
Under the six displayed input conditions, the construction
$T\mapsto N(T)$ supplies \emph{every} field of
\[
                       \mathrm{WN}(\mathcal D,a\circ z).
\]
In particular it proves the original
\texttt{Global\allowbreak Theory.\allowbreak Wach\allowbreak Lattice.\allowbreak build} implication, with its
original input hypotheses.

\begin{proof}
WLD18 supplies the actual $A$-submodule. WLD19 supplies
\texttt{N\_\allowbreak base}, \texttt{N\_\allowbreak fg} and \texttt{N\_\allowbreak proj};
WLD20 supplies \texttt{N\_\allowbreak phi} and \texttt{N\_\allowbreak phiinv};
WLD21 supplies \texttt{N\_\allowbreak Gamma} and \texttt{N\_\allowbreak pi};
WLD22 supplies \texttt{map\_\allowbreak N}; and WLD23 supplies
\texttt{one\_\allowbreak e}. By definition $N(T)=0$ for noncrystalline
$T$, giving \texttt{N\_\allowbreak junk}. These assertions use the same
$N(T)$, $\pi_A$, $\mathcal Q$ and the same given $D_T$ throughout.

The source's denominator really is invertible in $E$. Indeed
$\varphi_A(\pi_A)=\pi_Aq_0$, and both $\pi_A$ and
$\varphi_E(\pi_A)$ are units there. Therefore $q_0$ is a unit;
applying $\varphi_E$ another $s-1$ times proves that
$\mathcal Q=\varphi_A^{s-1}(q_0)$ is a unit. No additional
height bound or change of parameter is used.
\end{proof}
% END WLD24

% BEGIN WLD25
\paragraph{WLD25: the full prismatic lattice assembly.}
\label{claim:WLD25}
Assume the exact original
\texttt{Bhatt\allowbreak Scholze\allowbreak Crystalline\allowbreak Lattices} input, together with the
specified local prism, its maps, and the scalar data in
\texttt{Wach\allowbreak Lattice.\allowbreak build\_\allowbreak prism}. Thus the prism ring $S$ is local
and finite free over the specified $\mathbb Z_p[[X]]$; its
specialization has kernel $(\pi)$; its ideal is $(Q)$; its
embedding $a:S\hookrightarrow A$ has $a(\pi)=\pi_A$ and
$a(Q)=\mathcal Q$; its Galois maps commute with the distinguished
embedding and act trivially on the specialization. The supplied
basis of $O$ gives the stated $S$-coordinate decomposition of
$A$, the stated $W$-coordinates of $B$, and the tensor
identifications (WLD.1). The map $\iota$ is injective and reflects
units, $c$ reflects units, and the original fixed-tuple descent
hypothesis holds for the stated $H$ with trivial image in $\Gamma$.
For every crystalline $T$, also retain the original realization
datum \texttt{hreal}: a crystal and its Witt-tensor realization,
with its $\mathbb Z_p$-linear action equal to the specified action
on $T$. This datum is an explicit premise of the generic theorem;
the all-layer specialization will supply it from the published input.
Then the original output
\[
                         \mathrm{Nonempty}\bigl(
                         \mathrm{WN}(\mathcal D,a\circ z)\bigr)
                                                               \tag{WLD.26}
\]
is proved.

\begin{proof}
Choose a realizing crystal for each crystalline $T$, as allowed
by the supplied realization data. In the all-layer application
these data come from the existence clause of the crystalline-lattice
input. Keep that choice fixed for this $T$ and set
\[
             L_T=\kappa_T\bigl(j_T^{\mathrm{crys}}
                                  (\mathcal C_T(i)M_T)\bigr).
\]
Here $j_T^{\mathrm{crys}}$ is the crystalline realization, distinct
from the given \'etale embedding $j_T$; this notation makes the
two maps explicit.

By WLD24's last paragraph, $\mathcal Q$ is a unit in $E$.
The identity $\iota(a(Q))=c(\psi Q)$ and nonzero $B$
give $\psi(Q)\ne0$. Thus WLD11--WLD12 apply.
WLD1 transfers their equations to $Y_T$, and WLD17 gives
the descended basis and (WLD.20). The two Frobenius inclusions
and the group divisibility in item 4 are exactly (WLD.13)--(WLD.14)
on the image, with the same $\pi$ and $Q$.

WLD13 proves item 5 for \emph{all} the original crystalline
lattice maps, including the scalar endomorphisms needed for
$A$-stability. WLD16 shows that the specified trivial object
is crystalline. WLD14 gives precisely item 6 for its generator.
WLD4 supplies the exact coefficient decomposition in $E$;
its $W$-completeness is the $p$-adic completeness of the original
Witt ring. The remaining scalar conditions are the displayed
input data of this theorem. All six conditions are therefore
proved, and WLD24 gives (WLD.26).

For the source's arbitrary generic coefficient ring $O$, the given
basis and injective map $O\to LL$ supply the rank-one generator
used in WLD14; no completeness or discrete valuation hypothesis
on that generic $O$ has been added. Naturality WLD13 works by
restriction to $\mathbb Z_p$ and the supplied tensor maps.
Thus this is the full original \texttt{build\_\allowbreak prism} implication,
not merely its specialization to one conveniently chosen \'etale model.
\end{proof}
% END WLD25

\subsection{The remaining all-layer input, stated without extra coherence}

% BEGIN WLD26
\paragraph{WLD26: a precise three-supplier return to all-layer Wach existence.}
\label{claim:WLD26}
Fix an original prime $p$, local period datum $P$, tied integral
datum $A_{\inf}$, and retain
\texttt{Bhatt\allowbreak Scholze\allowbreak Crystalline\allowbreak Lattices}. Suppose the following
three source suppliers have been supplied for every original
$\mathrm{Good}(K,s,LL)$ and every permitted $\Gamma,R,I$:
\begin{enumerate}
 \item The actual Wach prism and its maps in
 \texttt{exists\_\allowbreak wach\_\allowbreak prism\_\allowbreak over\_\allowbreak integers}, including the
 scalar identity $a\circ z=\mathrm{zlA}$ and both finite
 coefficient-coordinate assertions.
 \item Injectivity and reflection of units for
 $\iota=I.\mathrm{ioE}$, the corresponding two fields of
 \texttt{wach\allowbreak Inclusion\_\allowbreak injective\_\allowbreak saturated\_\allowbreak inter}.
 \item The following homogeneous case of
 \texttt{fixed\allowbreak By\allowbreak Cyclotomic\allowbreak Kernel\_\allowbreak etale\allowbreak Phi\_\allowbreak solution\_\allowbreak mem\_\allowbreak wach}:
 for all $r,q$, $F\in\operatorname{GL}_r(E)$,
 $U\in\operatorname{GL}_q(E)$ and $V\in M_{r,q}(B)$ with
 cyclotomic-kernel-fixed entries,
 \[
                 \iota(F)\varphi_B(V)=V\iota(U)
                     \quad\Longrightarrow\quad V_{ij}\in\iota(E).
 \]
\end{enumerate}
Then the full original \texttt{Wach\allowbreak Ex} holds: for every such
$\Gamma,R,I$ and \emph{every supplied} \'etale module $\mathcal D$,
there is an $\mathrm{NS}(\mathcal D)$ with every original field.

\begin{proof}
Take any permitted layer and any of its allowed \'etale modules.
The first supplier gives exactly the local prism data used by
WLD25. Choose a finite $\mathbb Z_p$-basis of the ring of integers
$O$ of $LL$. The standard finite-free structure of the integers
in a finite $p$-adic field extension supplies this basis.
WLD1 supplies the original tensor cancellation maps; WLD2 supplies
the $W$-coordinates of $B=O\otimes_{\mathbb Z_p}W$ and reflection
of units by $c$. Its Frobenius and Galois compatibilities hold on
pure tensors, as checked in WLD1.

Let $H$ be the original cyclotomic kernel. Its image in $\Gamma$
is trivial: apply the injective character $\chi_\Gamma$, use
its specified cyclotomic-character equality, and then the
injectivity $O\hookrightarrow LL$.
WLD15, using suppliers 2 and 3, gives the fixed-tuple descent
required by WLD25. The existence clause of the crystalline-lattice
input provides the realizing crystal for every crystalline $T$;
its fullness clause is exactly the one used in WLD13--WLD14.
The tied local data give the required standard-period hypothesis.
WLD25 now yields \texttt{WN} with scalar map $a\circ z$.
The identity in supplier 1 identifies that map with the original
$\mathrm{zlA}$, so the output is literally \texttt{NS}.

All choices were made after the supplied layer and \'etale module
were fixed. The universal quantifiers over those data have not
been reduced. The existing DFD12 coefficient-extension argument
may supply the crystalline-lattice input from the original
integral Bhatt--Scholze input, but it does not supply any of the
three period-ring results retained above.
\end{proof}
% END WLD26

\paragraph{What this chapter supplies, and what remains.}
The completed original mathematical producers are the coefficient
lifting theorem, the denominator-matrix theorem, the base-change
lattice theorem, the finite-free restriction theorem, the local-prism
crystalline value theorem, and both generic Wach-lattice constructors.
They now have one detailed human argument with their actual common
objects. The matrix descent in WLD15 is proved from its explicitly
stated coefficient-descent premise; WLD26 is a conditional all-layer
return with exactly the three displayed period-ring suppliers.
Those suppliers remain mathematical tasks, not declarations proved
by this chapter.

The next useful period-ring target is the original injectivity and
unit-reflection theorem, or the homogeneous fixed-Frobenius matrix
descent at the exact scope in WLD26. The prism construction is
independent of those two and can proceed in parallel. After full
\texttt{Wach\allowbreak Ex}, the faithful-invariant and \'etale-basis work,
the Herr class maps, and the quantitative de Rham comparison are
still needed for the complete \texttt{HW/\allowbreak Fon/\allowbreak Cmp} triple.
The \texttt{Layers} assembly permits a whole complete record to be
chosen at each allowed outer index; it does not add an unspecified
inter-layer transition condition. SSC's completed signature comparison
and the later local family bounds retain their original scopes.
No whole \texttt{GLOBAL\_\allowbreak DENOMINATOR} closure is claimed.

\paragraph{Source correspondence.}
Every source below is at the accepted commit:
\par\noindent\texttt{7cbf7a0c32fca1badf05931879a37c3f8a2b03b4}.
\par\noindent
The exact source ranges and their hashes are recorded with each claim.
Under \texttt{Proofs/\allowbreak Inertia\allowbreak Eigenvalues}, the primary producers are
\texttt{Exists\allowbreak Witt\allowbreak Tilt\allowbreak Lattice\allowbreak Tensor\allowbreak Cancel},
\texttt{Exists\allowbreak Adic\allowbreak Coordinates\allowbreak Of\allowbreak Basis\allowbreak Mod\allowbreak P},
\texttt{Exists\allowbreak Is\allowbreak Unit\allowbreak Matrix\allowbreak Of\allowbreak Pow\allowbreak Smul\allowbreak Span\allowbreak Invert},
\texttt{Is\allowbreak Base\allowbreak Change\allowbreak Lattice\allowbreak Of\allowbreak Basis\allowbreak After\allowbreak Extension},
\texttt{Projective\allowbreak Of\allowbreak Free\allowbreak Lattice\allowbreak Over\allowbreak Finite\allowbreak Free\allowbreak Ring},
\texttt{Prismatic\allowbreak FCrystal\allowbreak Value\allowbreak Local\allowbreak Prism\allowbreak Ainf\allowbreak Witt\allowbreak Tensor},
\texttt{Witt\allowbreak Twist\allowbreak Invariant\allowbreak Basis\allowbreak Change},
\texttt{Is\allowbreak Crys\allowbreak Lat\allowbreak Of\allowbreak Trivial}, \texttt{Is\allowbreak Unit\allowbreak Of\allowbreak One\allowbreak Tmul}, and
\texttt{Wach\allowbreak Lattices\allowbreak Exist\allowbreak All\allowbreak Layers\allowbreak Of\allowbreak Bhatt\allowbreak Scholze}.
The unchanged \texttt{WN}, \texttt{WEt} and \texttt{Layers}
contracts are in \texttt{Definitions/\allowbreak Non\allowbreak Galois\allowbreak Descent/\allowbreak Exists\allowbreak Of\allowbreak Layers};
the all-layer quantifiers are in \texttt{Conv\allowbreak Zero}.





